\section{Problems}

\begin{pro}
How many license plates can be made using either three letters
followed by three digits or four letters followed by two digits?
\end{pro}
\begin{pro}
How many license plates can be made using 4 letters and 3 numbers
if the letters cannot be repeated and the letters and numbers may appear in any order?
\end{pro}
\begin{pro} 
How many bit strings of length 8 either begin with 
three 1s or end with four 0s?
\end{pro}
\begin{pro} 
How many alphabetic strings are there whose length is at most 5?
\end{pro}
\begin{pro}
How many bit strings are there of length at least 4 and at most 6?
\end{pro}
\begin{pro}
How many subsets with 4 or more elements does a set of size 30 have?
\end{pro}
% Now this is an evaluation example.
%\begin{pro} 
%Prove that at a gathering of $n\geq 2$ people, there are 
%two people who have shaken hands with the same number of
%people.
%\end{pro}
\begin{pro} 
Given a group of ten people, prove that at least 4 are male
or at least 7 are female.
\end{pro}
\begin{pro} 
My family wants to take a group picture.  There are 7 men
and 5 women, and we want none of the women to stand next to 
each other.  How many different ways are there for us to 
line up?
\end{pro}
\begin{pro} 
My family (7 men and 5 women) wants to select a group of
5 of us to plan Christmas.  We want at least
1 man and 1 woman in the group.  How many ways are
there for us to select the members of this group?
\end{pro}
\begin{pro} 
Compute each of the following:
$\binom{8}{4}$, $\binom{9}{9}$, $\binom{7}{3}$, $8!$, and $5!$
\end{pro}
\begin{pro}
For what value(s) of $k$ is $\binom{18}{k}$ largest?  smallest?
\end{pro}
\begin{pro}
For what value(s) of $k$ is $\binom{19}{k}$ largest?  smallest?
\end{pro}
\begin{pro}
A computer network consists of 10 computers.  
Each computer is directly connected to zero or more of the other computers.
\begin{lenum}
\item
Prove that there are at least two computers in the network that are directly connected to the
same number of other computers.
\item
Prove that there are an even number of computers that are connected to an odd number of
other computers.
\end{lenum}
\end{pro}
\begin{pro}
Simplify the following expression so it does not involve any factorials or binomial coefficients:
$\binom{x}{y}/\binom{x+1}{y-1}$
%${x \choose y}/{x+1\choose y-1}$.
\end{pro}

\begin{pro}
Prove that amongst six people in a room there are at least three who
know one another, or at least three who do not know one
another.\end{pro}

\begin{pro}
%[MMPC 1992] 
Suppose that the letters of the English
alphabet are listed in an arbitrary order. \begin{lenum} \item
Prove that there must be four consecutive consonants. \item Give a
list to show that there need not be five consecutive consonants.
\item Suppose that all the letters are arranged in a circle. Prove
that there must be five consecutive consonants.
\end{lenum}

\end{pro}
\begin{pro}
%[Stanford 1953] 
Bob has ten pockets and forty four silver
dollars. He wants to put his dollars into his pockets so
distributed that each pocket contains a different number of
dollars.
\begin{lenum}\item Can he do so? \item Generalize the problem, considering $p$ pockets
and $n$ dollars. Why is the problem most interesting when 
$n =
\dfrac{(p - 1)(p - 2)}{2}$? \end{lenum} \end{pro}

\begin{pro}
Expand and simplify the following.
\begin{lenum}
\item $(x - 4y)^3$ 
\item $(x^3 + y^2)^4$ 
\item $(2 + 3x)^6$ 
\item $(2i - 3)^5$ 
\item $(2i + 3)^4 + (2i - 3)^4$ 
\item $(2i + 3)^4 - (2i - 3)^4$ 
\item $(\sqrt{3} - \sqrt{2})^3$ 
\item $(\sqrt{3} +\sqrt{2})^3 + (\sqrt{3} - \sqrt{2})^3$ 
\item $(\sqrt{3} +\sqrt{2})^3 - (\sqrt{3} - \sqrt{2})^3$
\end{lenum}
\end{pro}

\begin{pro}
What is the coefficient of $x^6y^9$ in $(3x-2y)^{15}$?
\end{pro}
\begin{pro} What is the coefficient of $x^4 y^6$ in $(x\sqrt{2} - y)^{10}?$
\end{pro}

\begin{pro}\label{pro:pascident}
Prove Pascal's Identity (Theorem~\ref{thm:pascalIdent}).
(Hint: Just use the definition of the binomial coefficient and do a little algebra.)
\end{pro}

\begin{pro}
Prove that for any positive integer $n$, $\displaystyle \sum_{k=0}^n(-2)^k \binom{n}{k} = (-1)^n$.
(Hint: {\em Don't} use induction.)
\end{pro}

%\begin{pro}
%Prove that $$(a + b + c)^2 = a^2 + b^2 + c^2 + 2(ab + bc + ca)$$
%Prove that $$(a + b + c + d)^2 = a^2 + b^2 + c^2 + d^2 + 2(ab + ac +
%ad + bc + bd + cd )$$
%Generalize.
%\end{pro}
%
%\begin{pro}
%Compute $(x + 2y + 3z)^2$.
%\end{pro}

%\begin{pro}
%Given that $a + 2b = -8, \ ab = 4,$ find (i) $a^2 + 4b^2$, (ii) $a^3
%+ 8b^3$, (iii) $\dis{\frac{1}{a} + \frac{1}{2b}}$.
%\end{pro}

%\begin{pro}
%The sum of the squares of three consecutive positive integers is
%$21170$. Find the sum of the cubes of those three consecutive
%positive integers.
%\end{pro}

\begin{pro} Expand and simplify $$ (\sqrt{1 - x^2} + 1)^7 -
(\sqrt{1 - x^2} - 1)^7.$$\end{pro}
%%
%My problems.
\begin{pro}
There are approximately 7,000,000,000 people on the planet.
Assume that everyone has a name that consists of exactly $k$ 
lower-case letters from the English alphabet.
\begin{lenum}
\item
If $k=8$, is it guaranteed that two people have the same name?  Explain.
\item
What is the maximum value of $k$ that would guarantee that at least two people have the same name?
\item
What is the maximum value of $k$ that would guarantee that at least 100 people have the same name?
\item
Now assume that names can be between 1 and $k$ characters long.
What is the maximum value of $k$ that would guarantee that at least two people have the same name?
\end{lenum}
\end{pro}
\begin{pro}
Password cracking is the process of determining someone's password, typically using a computer.
One way to crack passwords is to perform an exhaustive search that tries every possible 
string of a given length until it (hopefully) finds it.
Assume your computer can test 10,000,000 passwords per second.
How long would it take to crack passwords with the following restrictions? Give answers in seconds,
minutes, hours, days, or years depending on how large the answer is (e.g. 12,344,440 seconds isn't 
very helpful).  Start by determining how many possible passwords there are in each case.
\begin{lenum}
\item
8 lower-case alphabetic characters.
\item
8 alphabetic characters (upper or lower).
\item
8 alphabetic (upper or lower) and numeric characters.
\item
8 alphabetic (upper or lower), numeric characters, and special characters 
(assume there are 32 allowable special characters).
\item
8 or fewer alphabetic (upper or lower) and numeric characters.
\item
10 alphabetic (upper or lower), numeric characters, and special characters 
(assume there are 32 allowable special characters).
\item
8 characters, with at least one upper-case, one lower-case, one number, and one special character. 
\end{lenum}
\end{pro}
\begin{pro}
IP addresses are used to identify computers on a network.  In IPv4, IP addresses are 32 bits long.
They are usually written using dotted-decimal notation, where the 32 bits are split up into 4 8-bit segments,
and each 8-bit segment is represented in decimal.  So the IP address $10000001$ $11000000$ $00011011$ $00000100$
is represented as $129.192.27.4$.
The {\em subnet mask} of a network is a string of $k$ ones followed by $32-k$ zeros, where the value of $k$
can be different on different networks.  For instance, the subnet mask might be 
$11111111111111111111111100000000$, which is $255.255.255.0$ in dotted decimal.
To determine the {\em netid}, an IP address is bitwise ANDed with the subnet mask.  
To determine the {\em hostid}, an IP address is bitwise ANDed with the bitwise complement of the subnet mask.
Since every computer on a network needs to have a different {\em hostid}, the number of possible 
{\em hostid}s determines the maximum number of computers that can be on a network.

Assume that the subnet mask on my computer is currently $255.255.255.0$ and my IP address is $209.140.209.27$.
\begin{lenum}
\item
What are the {\em netid} and {\em hostid} of my computer?
\item
How many computers can be on the network that my computer is on?
\item
In 2010, Hope College's network was not split into subnetworks like it is currently, so all of the
computers were on a single network that had a subnet mask of $255.255.240.0$.  
How many computers could be on Hope's network in 2010?
\end{lenum}
\end{pro}
\begin{pro}
Prove that ${\displaystyle \sum_{k=0}^{n} \binom{n}{k}=2^n}$ by counting the number of binary strings
of length $n$ in two ways.
\end{pro}

\begin{pro}
In March of every year people fill out brackets for the NCAA Basketball Tournament. They pick the
winner of each game in each round.  We will assume the tournament starts with 64 teams 
(it has become a little more complicated than this recently).
The first round of the tournament consists of 32 games, 
the second 16 games, the third 8, the fourth 4, the fifth 2, and the final 1.  
So the total number of games is $32+16+8+4+2+1=63$.
You can arrive at the number of games in a different way. Every game has a loser who is
out of the tournament.  Since only 1 of the 64 teams remains at the end, there must be 63 losers,
so there must be 63 games.
Notice that we can also write $1+2+4+8+16+32=63$ as ${\displaystyle\sum_{k=0}^5 2^k=2^6-1}$.
\begin{lenum}
\item 
Use a combinatorial proof to show that for any $n>0$, ${\displaystyle\sum_{k=0}^n 2^k=2^{n+1}-1}$.
(That is, define an appropriate set and count the cardinality of the set in two ways to
obtain the identity.)
\item
When you fill out a bracket you are picking who you think the winner will be of each game.  
How many different ways are there to fill out a bracket? 
(Hint: If you think about this in the proper way, this is pretty easy.)
\item
If everyone on the planet (7,000,000,000) filled out a bracket, is it guaranteed that 
two people will have the same bracket?  Explain.
\item
Assume that everyone on the planet fills out $k$ different brackets and that no brackets are
repeated (either by an individual or by anybody else).  How large would $k$ have to be before 
it is guaranteed that somebody has a bracket that correctly predicts the winner of every game?
\item
Assume every pair of people on the planet gets together to fill out a bracket (so everyone has 
6,999,999 brackets, one with every other person on the planet).  What is the smallest and largest 
number of possible repeated brackets?
\end{lenum}
\end{pro}
\begin{pro}
Mega Millions has 56 white balls, numbered 1 through 56, and 46 red balls, numbered 1 through 46.  
To play you pick 5 numbers between 1 and 56 (corresponding to white balls) and 1 number between 1 and 46
(corresponding to a red ball).  Then 5 of the 56 balls and 1 of the 46 balls are drawn
randomly (or so they would have us believe).  You win if your numbers match all 6 balls.
\begin{lenum}
\item
How many different draws are possible?
\item
If everyone in the U.S.A. bought a ticket (about 314,000,000), is it guaranteed that two people have the
same numbers?  Three people?
\item
If everyone in the U.S.A. bought a ticket, what is the maximum number of people 
that are guaranteed to share the jackpot?
\item
Which is more likely: Winning Mega Millions or picking every winner in the NCAA Basketball Tournament
(see previous question)?  How many more times likely is one than the other?
\item
I purchased a ticket last week and was surprised when {\em none} of my six numbers matched.  
Should I have been surprised?  
What are the chances that a randomly selected ticket will match none of the numbers?
\item
(hard) What is the largest value of $k$ such that you are more likely to pick at least $k$ winners in the 
NCAA Basketball Tournament than you are to win Mega Millions?
\end{lenum}
\end{pro}

\begin{pro}
You get a new job and your boss gives you 2 choices for your salary.  You can either make \$100 per day
or you can start at \$.01 on the first day and have your salary doubled every day.  You know that you will
work for $k$ days.  For what values of $k$ should you take the first offer and for
which should you take the second offer?  Explain.
\end{pro}

\begin{pro}\label{prob:csCourses}
The 300-level courses in the CS department are split into three groups: Foundations (361, 385), 
Applications (321, 342, 392), and Systems (335, 354, 376).  
In order to get a BS in computer science at Hope you need to take at least one course from
each group.
\begin{lenum}
\item
How many different ways are there of satisfying this requirement by taking exactly 3 courses?
\item
If you take four 300-level courses, 
how many different possibilities do you have that satisfy the requirements?
\item 
How many total ways are there to take 300-level courses that satisfy the requirements?
\item
What is the smallest $k$ such that no matter which $k$ 300-level courses you choose,
it is guaranteed that you will satisfy the requirement?
\end{lenum}
\end{pro}

\begin{pro}
I am implementing a data structure that consists of $k$ lists. 
I want to store a total of $n$ objects in this data structure, with each item being stored on one of the lists.
All of the lists will have the same capacity (e.g. perhaps each list can hold up to 10 elements).\\
Write a method \scode{minimumCapacity(int n, int k)} that computes the minimum capacity
each of the $k$ lists must have to accommodate $n$ objects.
In other words, if the capacity is less than this, then there is no way the objects can all be stored on the lists.
You may assume integer arithmetic truncates (essentially giving you the {\em floor} function), 
but that there is no {\em ceiling} function available.
\end{pro}

\begin{pro}
Write a method \verb+choose(int n, int k)+ (in a Java-like language) that computes $\binom{n}{ k}$.
Your implementation should be as efficient as possible.
Make sure to give and prove the efficiency of your algorithm.
\end{pro}


\begin{pro}
Let $A = \begin{bmatrix} 1 & 1 & 1 \cr
0 & 1 & 1 \cr 0 & 0 & 1 \cr
\end{bmatrix}. $
Prove that
$A^n = \begin{bmatrix} 1 & n & \frac{n(n + 1)}{2} \cr 0 & 1 & n \cr
0 & 0 & 1 \cr
\end{bmatrix}.$
%\begin{answer} 
%Put $$ J =    \begin{bmatrix} 0 & 1 & 1 \cr 0 & 0 & 1 \cr 0 & 0 & 0 \cr
%\end{bmatrix}. $$
%We first notice that
%$$J^2 = \begin{bmatrix} 0 & 0 & 1 \cr 0 &  0 & 0 \cr 0 &  0 &  0 \cr\end{bmatrix}, \ \ J^3 = {\bf 0}_3.$$
%This means that the sum in the binomial expansion $$A^n = ({\bf
%I}_3 + J)^n = \sum _{k = 0} ^n \binom{n}{k}{\bf I}^{n - k}J^k$$ is
%a sum of zero matrices for $k \geq 3.$ We thus have
%$$\begin{array}{lll}A^n & = & {\bf I}^n _3 + n{\bf I}^{n - 1} _3J + \binom{n}{2}{\bf I}^{n - 2}_3J^2 \vspace{2mm}\\
%&  = &
%\begin{bmatrix} 1 & 0 & 0 \cr 0 &  1 & 0 \cr 0 &  0 &  1 \cr\end{bmatrix} +
%\begin{bmatrix} 0 & n & n \cr 0 &  0 & n \cr 0 &  0 &  0 \cr\end{bmatrix}
%+ \begin{bmatrix} 0 & 0 & \binom{n}{2} \cr 0 &  0 & 0 \cr 0 &  0 &
%0 \cr\end{bmatrix} \vspace{2mm}\\
%&  = &  \begin{bmatrix} 1 & n & \frac{n(n + 1)}{2} \cr 0 & 1 & n
%\cr 0 & 0 & 1 \cr
%\end{bmatrix}, \end{array}$$giving the result, since $\binom{n}{2} = \frac{n(n - 1)}{2}$ and $n + \binom{n}{2} = \frac{n(n + 1)}{2}$.
%\end{answer}
\end{pro}